\documentclass[10pt,a4paper]{amsart}
\usepackage[margin=19mm]{geometry}
\usepackage{amsmath,amssymb,mathtools}
\usepackage[scr=rsfs,cal=euler]{mathalfa}
\usepackage{xfrac}
\usepackage[unicode,hidelinks,bookmarks=false]{hyperref}
\newcommand{\ie}{i.e.\ }
\newcommand{\cf}{cf.\ }
\newcommand{\resp}{respectively\ }
\newcommand{\Subsection}{\subsection}
\providecommand{\nobreakdash}{\nobreak\hbox{-}}
\setlength{\emergencystretch}{2em}
\multlinegap0pt
\usepackage{bm}
\usepackage{bbm}
\usepackage{dsfont}
\usepackage{xspace}
\usepackage{cancel}
\usepackage{mathtools}
\usepackage{amsthm}
\makeatletter
\@ifundefined{theorem}{\newtheorem{theorem}{Theorem}}{}
\@ifundefined{lemma}{\newtheorem{lemma}[theorem]{Lemma}}{}
\makeatother
\let\div\relax\DeclareMathOperator{\div}{div}
\newcommand{\R}{\mathbb{R}}
\newcommand{\W}{\mathcal{W}}
\title[Variational problems in $L^\infty$ involving semilinear seco]{Variational problems in $L^\infty$ involving semilinear second order differential operators}
\author{Nikos Katzourakis, Roger Moser}
\date{Source: arXiv:2303.15982v3 (CC BY 4.0)}
\begin{document}
\maketitle

\noindent\emph{Source and licence.} Variational problems in $L^\infty$ involving semilinear second order differential operators. Version arXiv:2303.15982v3 (CC BY 4.0), distributed under the Creative Commons Attribution 4.0 International licence, as recorded in the arXiv licence metadata for this exact version and shown on the abstract page https://arxiv.org/abs/2303.15982v3. Full rights notice: licenses/arxiv-2303.15982-CC-BY-4.0.txt.\par\smallskip

\noindent\textit{Editorial scope.} Counted unit: the Lemma whose source label is \texttt{lem:distributional-solutions} in the article (source0.tex lines 520--527, source label \texttt{lem:distributional-solutions}), delivered together with the article's own complete original proof of it (source0.tex lines 529--622, one \texttt{proof} environment of 94 lines). No mathematics is written, repaired or re-derived: statement and proof are the article's own text line for line. Required context retained uncounted: eq:div-div-homogeneous (lines 515--517). Every definition, display and label of those lines is retained unchanged. Omitted from this package and not redistributed: 1--514, 518--519, 528--528, 623--1357. Nothing inside a retained excerpt is omitted or altered: every retained body is delivered exactly as the article's lines are written, so no omission receipt of any kind accompanies this package. Citations: the retained excerpts carry no citation command at all, so no bibliography entry of the article is transcribed and no reference is redistributed. Reference adaptation: none was needed and none was made; every reference inside the delivered unit resolves inside it, so no unresolved reference or citation is produced. Printed slips kept literal: no printed word, symbol, hypothesis or number of the counted statement or of its proof was altered. Layout and class: the article's own class is \texttt{article} with packages including amsfonts, latexsym, amssymb, amsthm, amsmath, mathrsfs, esint, authblk, hyperref; the wrapper is the lane's pinned \texttt{amsart} editorial wrapper at 10pt on A4, which restates the theorem-like environments the retained text needs and adds the article's own macro definitions that the retained text needs (\texttt{\textbackslash R}, \texttt{\textbackslash W}, \texttt{\textbackslash div}), with extra wrapper packages (bm, bbm, dsfont, xspace, cancel, mathtools). The delivered numbering is the unit's own. Assets: no retained span contains a figure environment, a graphics-inclusion command, a PDF-page inclusion, a table or a TikZ picture; no image file is copied into this package, the package declares no asset dependency and no image file is redistributed with it. Licence: version arXiv:2303.15982v3 of this article is distributed under the Creative Commons Attribution 4.0 International licence, as recorded in the arXiv licence metadata for this exact version and shown on the abstract page https://arxiv.org/abs/2303.15982v3. A full rights notice is delivered at licenses/arxiv-2303.15982-CC-BY-4.0.txt. Characters: every retained excerpt is pure ASCII, so the delivered engine, XeLaTeX with its default Latin Modern text font, reports no missing character.
\begin{equation} \label{eq:div-div-homogeneous}
\div \div (fA) + \div(fB) + cf = 0
\end{equation}

\begin{lemma} \label{lem:distributional-solutions}
Suppose that $f \in L^1(\Omega)$ is a weak solution of equation
\eqref{eq:div-div-homogeneous}. Then
\[
\int_\Omega f(A : D^2\phi - B \cdot D\phi + c\phi) \, dx = 0
\]
for all $\phi \in \W_0^{2, \infty}(\Omega)$.
\end{lemma}

\begin{proof}
Given $\phi \in \W_0^{2, \infty}(\Omega)$, we first construct a family of
approximations $(\phi_\epsilon)_{\epsilon \in (0, \epsilon_0]}$ in
$C_0^\infty(\Omega)$ such that $\phi_\epsilon \to \phi$ in $W^{2, q}(\Omega)$
for any $q < \infty$ and, at the same time, such that $A : D^2\phi_\epsilon$
remains bounded in $L^\infty(\Omega)$.

For this purpose, we extend $\phi$ by $0$ outside of $\Omega$.
Choose a finite open cover $\{G_1, \dotsc, G_L\}$ of
$\overline{\Omega}$ with the property that there exist $R > 0$ and
there exist open cones $C_1, \dotsc, C_L \subseteq \R^n$ such that
for any $x \in \partial \Omega \cap G_\ell$,
\[
C_\ell \cap B_R(0) \cap (x - \overline{\Omega}) = \emptyset.
\]
(Hence $(x - C_\ell) \cap B_R(x)$ is an exterior cone to $\overline{\Omega}$.)
This is possible, because $\Omega$ is a bounded Lipschitz domain.

For every $\ell = 1, \dotsc L$, choose $\eta_\ell \in C_0^\infty(C_\ell \cap B_1(0))$ with $\eta_\ell \ge 0$ and
\[
\int_{B_1(0)} \eta_\ell(x) \, dx = 1.
\]
For $\epsilon \in (0, R]$, set
\[
\eta_{\ell\epsilon}(x) = \frac{1}{\epsilon^n} \eta_\ell\left(\frac{x}{\epsilon}\right)
\]
and
\[
\phi_{\ell\epsilon} = \phi * \eta_{\ell\epsilon}.
\]
Then for $x \in \partial \Omega \cap G_\ell$,
\[
\phi_{\ell\epsilon}(x) = \int_{C_\ell \cap B_\epsilon(0) \cap (x - \Omega)} \eta_{\ell\epsilon}(y) \phi(x - y) \, dy = 0.
\]
Moreover, the function $\phi_{\ell\epsilon}$ vanishes in a neighbourhood of
any such point $x$.

Now choose a partition of unity $\chi_1, \dotsc, \chi_L$ in $\Omega$ with
$\chi_\ell \in C_0^\infty(G_\ell)$ for $\ell = 1, \dotsc, L$. Set
\[
\phi_\epsilon = \sum_{\ell = 1}^L \chi_\ell \phi_{\ell\epsilon}.
\]
Then it is clear that $\phi_\epsilon \in C_0^\infty(\Omega)$ and that
$\phi_\epsilon \to \phi$ in $W^{2, q}(\Omega)$, for any $q < \infty$, as
$\epsilon \searrow 0$.

For $x \in \Omega$, we compute
\[
\begin{split}
A(x) : D^2 \phi_{\ell\epsilon}(x) &= \int_{B_\epsilon(0)} \eta_{\ell\epsilon}(y) A(x) : D^2\phi(x - y) \, dy \\
& = \int_{B_\epsilon(0)} \eta_{\ell\epsilon}(y) (A(x) - A(x - y)) : D^2\phi(x - y) \, dy \\
& \quad + \int_{B_\epsilon(0)} \eta_{\ell\epsilon}(y) A(x - y) : D^2\phi(x - y) \, dy.
\end{split}
\]
Using an integration by parts, we find that
\begin{multline*}
\int_{B_\epsilon(0)} \eta_{\ell\epsilon}(y) (A(x) - A(x - y)) : D^2\phi(x - y) \, dy \\
\begin{aligned}
&= \int_{B_\epsilon(0)} \eta_{\ell\epsilon}(y) \div A(x - y) \cdot D\phi(x - y) \, dy \\
& \quad + \int_{B_\epsilon(0)} (A(x) - A(x - y)) : D\eta_{\ell\epsilon}(y) \otimes D\phi(x - y) \, dy.
\end{aligned}
\end{multline*}
For $y \in B_\epsilon(0)$, we have the inequality
\[
|A(x) - A(x - y)| \le \epsilon \|A\|_{C^1(\overline{\Omega})}.
\]
Hence there exists a universal constant $C_1$ such that
\begin{multline*}
\left|\int_{B_\epsilon(0)} \eta_{\ell\epsilon}(y) (A(x) - A(x - y)) : D^2\phi(x - y) \, dy\right| \\
\le C_1\|A\|_{C^1(\overline{\Omega})} \bigl(1 + \|D\eta_\ell\|_{L^1(B_1(0))}\bigr) \|D\phi\|_{L^\infty(\Omega)}.
\end{multline*}
It is clear that
\[
\left|\int_{B_\epsilon(0)} \eta_{\ell\epsilon}(y) A(x - y) : D^2\phi(x - y) \, dy\right| \le \|A : D^2 \phi\|_{L^\infty(\Omega)}.
\]
Thus we obtain a uniform estimate for $\|A : D^2\phi_{\ell\epsilon}\|_{L^\infty(\Omega)}$. A similar estimate for $\phi_\epsilon$ is then easy to prove.

It follows that there exists a sequence $\epsilon_k \searrow 0$ such that
$A : D^2\phi_{\epsilon_k} \stackrel{*}{\rightharpoonup} g$, weakly* in
$L^\infty(\Omega)$, for some $g \in L^\infty(\Omega)$. For any $\psi \in C_0^\infty(\Omega)$, we then compute
\[
\begin{split}
\int_\Omega \psi g \, dx &= \lim_{k \to \infty} \int_\Omega \psi A : D^2 \phi_{\epsilon_k} \, dx \\
&= \lim_{k \to \infty} \int_\Omega \div \div (\psi A) \phi_{\epsilon_k} \, dx \\
&= \int_\Omega \div \div (\psi A) \phi \, dx \\
&= \int_\Omega \psi A : D^2 \phi \, dx.
\end{split}
\]
It follows that $g = A : D^2\phi$. It then also follows that $A : D^2\phi_\epsilon \stackrel{*}{\rightharpoonup} A : D^2\phi$, i.e, it
is not necessary to take a subsequence.

Now the claim of the lemma is proved by approximation with $\phi_\epsilon$
and with standard arguments.
\end{proof}

\end{document}
